% =============================================================================
% CHƯƠNG 3 — PHƯƠNG PHÁP NGHIÊN CỨU
% Văn phong khoa học, thuật ngữ và cấu trúc đã được chuẩn hóa.
% =============================================================================

\chapter{PHƯƠNG PHÁP NGHIÊN CỨU}
\label{chap:phuong-phap}

\section{Thiết kế nghiên cứu}
\label{sec:c3-design}

Nghiên cứu sử dụng phương pháp định lượng và thực nghiệm tính toán. Quy trình
gồm năm bước:
\begin{enumerate}[label=(\arabic*)]
    \item thu thập, căn chỉnh và chuẩn hóa giá ba tài sản cơ sở;
    \item hiệu chỉnh ba mô hình sinh dữ liệu trên cùng cửa sổ lịch sử và sinh các tập quỹ
    đạo có cùng kích thước;
    \item huấn luyện cùng một chính sách \textit{Deep Hedging} trên từng tập
    quỹ đạo bằng hai giai đoạn MSE và CVaR;
    \item đánh giá các checkpoint trên các cửa sổ lịch sử của ba giai đoạn thị
    trường;
    \item so sánh bằng chỉ tiêu mô tả, kiểm định ghép cặp có điều chỉnh nhiều
    phép thử và Model Confidence Set.
\end{enumerate}

Đơn vị thực nghiệm là một tổ hợp
$(\text{mô hình sinh dữ liệu},\text{hợp đồng},\kappa,\text{seed})$. Mỗi tổ hợp
tạo một checkpoint sau MSE và một checkpoint sau CVaR. Mỗi checkpoint được
đánh giá trên các cửa sổ của giai đoạn căng thẳng, giai đoạn phục hồi và giai
đoạn đại diện mà không huấn luyện lại mạng.

\section{Dữ liệu và xử lý dữ liệu}
\label{sec:c3-data}

\subsection{Nguồn dữ liệu và tài sản cơ sở}

Dữ liệu gồm giá đóng cửa điều chỉnh của ba cổ phiếu Hoa Kỳ: Apple (AAPL),
JPMorgan Chase (JPM) và ExxonMobil (XOM), được thu thập qua Yahoo Finance. Ba
tài sản thuộc các ngành công nghệ, tài chính và năng lượng. Các chuỗi được căn
chỉnh theo ngày giao dịch chung; giá trị thiếu sau căn chỉnh được điền bằng
quan sát gần nhất.

Lợi suất log của tài sản $i$ được tính bằng
\begin{equation}
    r_t^i=\log\left(\frac{S_t^i}{S_{t-1}^i}\right).
    \label{eq:c3-log-return}
\end{equation}
Trước khi đưa vào mạng hedging, mỗi quỹ đạo giá được chuẩn hóa theo giá đầu
kỳ:
\begin{equation}
    \widetilde S_t^i=\frac{S_t^i}{S_0^i},
    \qquad \widetilde S_0^i=1.
    \label{eq:c3-normalized-price}
\end{equation}
Sau chuẩn hóa, giá của mỗi tài sản bằng 1 tại đầu kỳ. Vì vậy, mức giá thực hiện
$\kappa$ được diễn giải theo tỷ lệ so với giá đầu kỳ và không phụ thuộc vào
mệnh giá riêng của từng cổ phiếu.

\subsection{Giai đoạn hiệu chỉnh và các giai đoạn đánh giá}

\begin{table}[H]
\centering
\small
\caption{Cấu trúc thời gian của dữ liệu; nhãn trong ngoặc là tên giai đoạn trong mã nguồn}
\label{tab:c3-periods}
\begin{tabular}{@{}llrl@{}}
\toprule
\textbf{Giai đoạn} & \textbf{Khoảng thời gian} & \textbf{Số ngày} &
\textbf{Vai trò}\\
\midrule
Hiệu chỉnh & 01/2005--12/2019 & 3.774 & Hiệu chỉnh ba mô hình sinh dữ liệu\\
Căng thẳng (COVID) & 01/2019--12/2020 & 503 & Đánh giá\\
Phục hồi (PostCOVID) & 01/2021--12/2022 & 503 & Đánh giá\\
Đại diện (Recent) & 01/2023--12/2025 & 750 & Đánh giá\\
\bottomrule
\end{tabular}
\end{table}

Giai đoạn căng thẳng bắt đầu từ năm 2019 để các cửa sổ trượt 252 ngày bao phủ cú sốc
tháng 2--4/2020. Năm 2019 đồng thời thuộc giai đoạn hiệu chỉnh mô hình sinh dữ
liệu. Do đó, giai đoạn căng thẳng không phải tập ngoài mẫu hoàn toàn đối với bước
hiệu chỉnh. Mạng hedging chỉ được huấn luyện trên quỹ đạo tổng hợp và không
được huấn luyện trực tiếp trên các cửa sổ giá lịch sử.

Từ 3.774 quan sát của giai đoạn hiệu chỉnh, các cửa sổ chồng lấn có độ dài
$N=252$ được xây dựng. Đối với SBTS, tập tham chiếu gồm
$M=3.522$ cửa sổ. Đối với đánh giá, mỗi giai đoạn cũng được chuyển thành tập các
cửa sổ chuẩn hóa dài 252 ngày theo cùng quy tắc.

\section{Hiệu chỉnh và sinh quỹ đạo}
\label{sec:c3-calibration}

Mỗi mô hình sinh dữ liệu tạo $L=20.000$ quỹ đạo gồm 252 ngày và ba tài sản.
Các quỹ đạo được chia thành 16.000 đường huấn luyện, 2.000 đường xác thực và
2.000 đường kiểm tra, với seed chia tập bằng 42. Kích thước ba tập được giữ cố
định giữa GBM, Heston và SBTS.

\subsection{Mô hình sinh dữ liệu GBM đa biến}

GBM được hiệu chỉnh bằng mô-men mẫu. Với $\Delta t=1/252$, bước sinh lợi suất
log là
\begin{equation}
    r_t^i=\left(\mu_i-\frac12\sigma_i^2\right)\Delta t
    +\sigma_i\sqrt{\Delta t}\,Z_t^i,
    \label{eq:c3-gbm-discrete}
\end{equation}
trong đó vector $Z_t$ được tạo bằng cách nhân nhiễu chuẩn độc lập với nhân
Cholesky của ma trận tương quan lợi suất thực nghiệm. Độ biến động năm ước
lượng cho AAPL, JPM và XOM lần lượt là $0{,}3232$, $0{,}3729$ và $0{,}2341$;
các tương quan cặp là $0{,}399$, $0{,}377$ và $0{,}468$.

\subsection{Mô hình sinh dữ liệu Heston đa biến}

Mỗi tài sản có một phương trình giá và một phương trình phương sai theo
\eqref{eq:c2-heston}. Đặc tả giản lược dùng chung
\begin{equation}
    \kappa_v=5{,}0,\qquad \xi=0{,}7,\qquad \rho=-0{,}7,
    \label{eq:c3-heston-shared}
\end{equation}
trong khi phương sai dài hạn và phương sai ban đầu của từng tài sản bằng bình
phương độ biến động lịch sử. Các cú sốc giá giữa tài sản được tương quan hóa
bằng cùng ma trận dùng cho GBM; cú sốc phương sai được tạo riêng theo từng tài
sản và ghép với cú sốc giá thông qua $\rho$. Mô phỏng dùng bốn bước nhỏ trong
mỗi ngày, full truncation để giữ phương sai không âm, Milstein cho phương sai
và log-Euler cho giá.

\subsection{Mô hình sinh dữ liệu SBTS đa biến}

SBTS sử dụng $M=3.522$ cửa sổ tham chiếu và kernel đa biến
\eqref{eq:c2-kernel}. Tại mỗi ngày mô phỏng, lịch sử $K$ ngày gần nhất của quỹ
đạo đang sinh được so sánh với các cửa sổ tham chiếu. Trọng số tương tự lịch
sử được kết hợp với nhân Brown suy ra từ nghiệm Cole--Hopf; chuyển tiếp là
trung bình có trọng số của các dịch chuyển kế tiếp trong mẫu, cộng với nhiễu
Brown ở lưới Euler. Cấu hình dùng $N^{\pi}=100$ bước nhỏ trong mỗi ngày.

Băng thông và bậc nhớ được chọn theo MSE trên tập giữ lại với 84 cấu hình. Cặp
tham số được sử dụng là
\begin{equation}
    (h^{*},K^{*})=(0{,}030,5),
    \qquad \mathrm{MSE}^{*}_{\mathrm{holdout}}=1{,}75\times10^{-4}.
    \label{eq:c3-hk-result}
\end{equation}
$K^{*}=5$ tương ứng với năm ngày giao dịch gần nhất.

\section{Hợp đồng quyền chọn được nghiên cứu}
\label{sec:c3-contracts}

Với $d=3$, $N=252$ và giá ban đầu chuẩn hóa bằng 1, giá trung bình theo thời
gian của tài sản $i$ là
\begin{equation}
    \bar R_i=\frac{1}{252}\sum_{n=1}^{252}\frac{S_n^i}{S_0^i}.
    \label{eq:c3-average-normalized-price}
\end{equation}
Khoản thanh toán của hai hợp đồng được xác định bởi
\begin{align}
    \Phi_{\mathrm{basket}}
    &=\left(\frac13\sum_{i=1}^{3}\bar R_i-\kappa\right)^{+},
    \label{eq:c3-basket}\\
    \Phi_{\mathrm{worst}}
    &=\left(\kappa-\min_{1\le i\le3}\bar R_i\right)^{+}.
    \label{eq:c3-worst}
\end{align}
Đối với basket Asian call, đại lượng
$\frac{1}{3}\sum_{i=1}^{3}\bar R_i$ là giá trung bình của rổ tài sản. Quyền
chọn có khoản thanh toán dương khi giá trung bình của rổ lớn hơn $\kappa$.
Đối với Asian worst-of put, $\min_i\bar R_i$ là giá trung bình của tài sản có
kết quả thấp nhất trong rổ. Quyền chọn có khoản thanh toán dương khi đại lượng
này thấp hơn $\kappa$.

Ba mức giá thực hiện là
$\kappa\in\{0{,}95;1{,}00;1{,}05\}$. Vì giá đầu kỳ đã được chuẩn hóa bằng 1,
$\kappa=0{,}95$, $1{,}00$ và $1{,}05$ lần lượt tương ứng với ngưỡng bằng 95\%,
100\% và 105\% giá đầu kỳ. Trạng thái trong giá, ngang giá hoặc ngoài giá được
xác định riêng cho quyền chọn mua và quyền chọn bán.

\section{Mô hình \textit{Deep Hedging}}
\label{sec:c3-hedging-model}

\subsection{Biến trạng thái và kiến trúc mạng}

Tại ngày $t$, vector trạng thái gồm
\begin{equation}
    X_t=\left(S_t^1,S_t^2,S_t^3,
    \bar S_t^1,\bar S_t^2,\bar S_t^3,
    \delta_{t-1}^1,\delta_{t-1}^2,\delta_{t-1}^3,T-t\right)
    \in\mathbb{R}^{10},
    \label{eq:c3-state}
\end{equation}
trong đó $\bar S_t^i$ là trung bình chạy và $\delta_{-1}=0$. Mạng truyền thẳng
có kiến trúc
\begin{equation}
    10\longrightarrow64\longrightarrow64\longrightarrow3,
    \label{eq:c3-architecture}
\end{equation}
với ReLU ở hai lớp ẩn và đầu ra tuyến tính là ba vị thế hedging. Trọng số
được khởi tạo bằng Xavier-uniform; hệ số khởi tạo của lớp cuối bằng $0{,}1$.
Mạng có 5.059 tham số. Khi bổ sung tham số vô hướng $V_0$, giai đoạn MSE có
5.060 tham số cần ước lượng.

\subsection{Giá trị danh mục, chi phí và sai số hedging}

Với chi phí giao dịch tỷ lệ $c=0{,}001$, giá trị cuối kỳ của danh mục hedging là
\begin{equation}
    V_T=V_0+\sum_{t=0}^{N-1}\delta_t^{\top}(S_{t+1}-S_t)
    -c\sum_{t=0}^{N-1}\sum_{i=1}^{3}
    |\delta_t^i-\delta_{t-1}^i|S_t^i.
    \label{eq:c3-terminal-wealth}
\end{equation}
Sai số hedging được định nghĩa theo nghĩa vụ của người bán:
\begin{equation}
    R=\Phi-V_T
    =\Phi-V_0-\mathrm{PnL}(\delta)+\mathrm{cost}(\delta).
    \label{eq:c3-residual}
\end{equation}
$R>0$ biểu thị thiếu hụt; do đó phần đuôi trên của $R$ là vùng tổn thất cần
được CVaR kiểm soát.

Trên mini-batch kích thước $B$, hai hàm mất mát là
\begin{align}
    \mathcal{L}_{\mathrm{MSE}}
    &=\frac{1}{B}\sum_{b=1}^{B}(R^{(b)})^2,
    \label{eq:c3-mse-loss}\\
    \mathcal{L}_{\mathrm{CVaR}}
    &=\nu+\frac{1}{(1-0{,}95)B}
      \sum_{b=1}^{B}(R^{(b)}-\nu)^{+}.
    \label{eq:c3-cvar-loss}
\end{align}

\section{Quy trình huấn luyện hai giai đoạn}
\label{sec:c3-curriculum}

\begin{table}[H]
\centering
\small
\caption{Cấu hình huấn luyện MSE--CVaR dùng chung cho toàn bộ lưới}
\label{tab:c3-hyperparameters}
\begin{tabular}{@{}lll@{}}
\toprule
\textbf{Tham số} & \textbf{Giai đoạn 1} & \textbf{Giai đoạn 2}\\
\midrule
Hàm mất mát & MSE & $\mathrm{CVaR}_{0{,}95}$\\
Tham số học & $\theta,V_0$ & $\theta,\nu$\\
Tham số cố định & Không & $V_0=V_0^{*}$\\
Learning rate & $10^{-3}$ & $10^{-4}$\\
Epoch tối đa & 500 & 200\\
Early-stopping patience & 20 & 20\\
Scheduler patience/factor & $10/0{,}5$ & $10/0{,}5$\\
Learning rate tối thiểu & $10^{-7}$ & $10^{-7}$\\
\midrule
Optimizer & \multicolumn{2}{c}{Adam, $\beta_1=0{,}9$, $\beta_2=0{,}999$}\\
Mini-batch & \multicolumn{2}{c}{$B=4.096$}\\
Gradient clipping & \multicolumn{2}{c}{Chuẩn $1{,}0$}\\
Chi phí giao dịch & \multicolumn{2}{c}{$c=0{,}001$}\\
\bottomrule
\end{tabular}
\end{table}

Giai đoạn 1 tối thiểu hóa MSE và ước lượng đồng thời $V_0$. Checkpoint có hàm
mất mát xác thực thấp nhất được lưu; giá trị tương ứng được ký hiệu là
$V_0^{*}$. Giai đoạn 2 tải checkpoint này, cố định $V_0^{*}$, khởi tạo $\nu$
từ phân vị 95\% của sai số hedging trên tập xác thực và tối thiểu hóa CVaR.
Adam, cắt gradient, dừng sớm và lịch giảm tốc độ học được áp dụng giống nhau
cho ba mô hình sinh dữ liệu \parencite{KingmaBa2015}.

MSE sử dụng toàn bộ quan sát trong mini-batch để khởi tạo tham số. CVaR chỉ sử
dụng các quan sát thuộc phần đuôi bất lợi. $V_0$ được cố định trong giai đoạn
CVaR vì tính tịnh tiến của CVaR làm bài toán tối ưu theo $V_0$ không có nghiệm
dừng hữu hạn.

\section{Thiết kế thực nghiệm và các yếu tố kiểm soát}
\label{sec:c3-experiment}

Lưới thực nghiệm là
\begin{equation}
    \underbrace{3}_{\text{mô hình sinh dữ liệu}}
    \times\underbrace{2}_{\text{hợp đồng}}
    \times\underbrace{3}_{\text{mức giá thực hiện}}
    \times\underbrace{10}_{\text{seed}}
    =180\ \text{cấu hình huấn luyện được giữ lại}.
    \label{eq:c3-grid}
\end{equation}
Mỗi cấu hình có hai checkpoint, tương ứng với 360 checkpoint. Mười seed được
đánh số từ 0 đến 9 và được dùng đồng bộ giữa các nhánh. Riêng cấu hình
SBTS--worst-of put--$\kappa=0{,}95$--seed 3 không hội tụ ở giai đoạn CVaR và
được thay bằng một lần chạy lại hội tụ.

\begin{table}[H]
\centering
\small
\caption{Các yếu tố được giữ cố định giữa ba mô hình sinh dữ liệu}
\label{tab:c3-fairness}
\begin{tabularx}{\textwidth}{@{}lX@{}}
\toprule
\textbf{Yếu tố} & \textbf{Giá trị chung}\\
\midrule
Cửa sổ hiệu chỉnh & 2005--2019, 3.774 ngày giao dịch\\
Số quỹ đạo & 20.000; chia 16.000/2.000/2.000\\
Độ dài và số tài sản & 252 ngày, 3 tài sản, $S_0^i=1$\\
Hợp đồng & Basket Asian call và Asian worst-of put\\
Mức giá thực hiện & $0{,}95;1{,}00;1{,}05$\\
Kiến trúc & $10\to64\to64\to3$, ReLU\\
Huấn luyện & MSE tối đa 500 epoch, CVaR tối đa 200 epoch\\
Chi phí giao dịch & $c=0{,}001$\\
Đánh giá & Cùng các cửa sổ của ba giai đoạn thị trường\\
\bottomrule
\end{tabularx}
\end{table}

Các chỉ tiêu gồm P\&L trung bình, độ lệch chuẩn của sai số hedging
$\sigma(R)$, $\mathrm{CVaR}_{0{,}95}(R)$, khoản thanh toán, chi phí giao dịch và
$V_0^{*}$. P\&L được xác định bằng $V_T-\Phi=-R$. Giá trị P\&L trung bình gần
0 biểu thị sai số trung bình nhỏ; $\sigma(R)$ thấp biểu thị kết quả lãi/lỗ ổn
định hơn; $\mathrm{CVaR}_{0{,}95}(R)$ thấp biểu thị mức thiếu hụt trung bình nhỏ
hơn trong 5\% trường hợp bất lợi nhất. Chỉ tiêu xếp hạng chính trong Chương 4
là $\sigma(R)$ tại checkpoint sau CVaR.

\section{Phương pháp suy luận thống kê}
\label{sec:c3-statistics}

\subsection{Kiểm định $t$ ghép cặp và kích thước ảnh hưởng}

Với hai mô hình sinh dữ liệu $A$ và $B$, đặt $d_s=L_s^A-L_s^B$, trong đó
$L_s$ là chỉ tiêu đánh giá ở seed $s$. Thống kê kiểm định là
\begin{equation}
    t=\frac{\bar d}{s_d/\sqrt{n}},\qquad n=10,
    \label{eq:c3-paired-t}
\end{equation}
và Cohen's $d$ ghép cặp là $d_{\mathrm{Cohen}}=\bar d/s_d$. Ghép seed kiểm
soát biến thiên do khởi tạo mạng.

\subsection{Điều chỉnh Benjamini--Hochberg}

Thiết kế tạo 216 phép kiểm định:
\begin{equation}
\begin{aligned}
216={}&(3\ \text{cặp})\times(2\ \text{hợp đồng})
\times(3\ \text{mức giá thực hiện})\\
&\times(3\ \text{giai đoạn})
\times(2\ \text{giai đoạn huấn luyện})\\
&\times(2\ \text{chỉ tiêu}).
\end{aligned}
\end{equation}
Toàn bộ $p$-value được điều chỉnh chung bằng Benjamini--Hochberg ở mức
$q=0{,}05$ \parencite{BenjaminiHochberg1995}. Một mô hình được xếp hạng cao
hơn khi có chỉ tiêu trung bình thấp hơn và $p_{\mathrm{BH}}<0{,}05$. Các
trường hợp còn lại được ghi nhận là chưa có khác biệt thống kê.

\subsection{Model Confidence Set}

MCS của \textcite{HansenLundeNason2011} được áp dụng riêng cho từng tổ hợp
hợp đồng--mức giá thực hiện--giai đoạn trên $\sigma(R)$ của checkpoint CVaR. Thủ tục
dùng mức ý nghĩa $\alpha=0{,}10$ và 5.000 bootstrap ở cấp seed. MCS loại dần
mô hình có chỉ tiêu đánh giá cao nhất cho đến khi không còn bác bỏ giả thuyết các mô hình còn
lại có cùng giá trị kỳ vọng của chỉ tiêu. Tập một phần tử biểu thị một mô hình sinh dữ liệu được lựa chọn rõ; tập nhiều
phần tử biểu thị dữ liệu chưa đủ để phân biệt các ứng viên.

Kiểm định ghép cặp sử dụng chênh lệch trong từng seed. MCS đánh giá đồng thời
ba mô hình và sử dụng độ phân tán giữa các seed. Hai phương pháp có thể tạo tập
kết quả khác nhau vì sử dụng cấu trúc thông tin khác nhau.

\section{Quy trình đối chiếu kết quả và khả năng tái lập}
\label{sec:c3-reproducibility}

Nguồn kết quả gồm ba notebook:
\begin{itemize}
    \item \texttt{article\_gbm.ipynb};
    \item \texttt{article\_heston\_sbts.ipynb};
    \item \texttt{article\_4\_3period.ipynb}.
\end{itemize}
Hai notebook đầu tạo dữ liệu, huấn luyện và lưu checkpoint. Notebook thứ ba
đánh giá checkpoint trên ba giai đoạn, tạo bảng kiểm định, MCS, $V_0^{*}$ và
kết quả của hai giai đoạn huấn luyện. Các notebook có tiền tố \texttt{25\_}
chỉ được dùng để đối chiếu lịch sử và không được dùng làm nguồn kết quả cuối.

Các số liệu tại Chương 4 được lấy từ bảng đầu ra của notebook tổng hợp. Các
phân tích không có trong đầu ra cuối cùng không được sử dụng làm bằng chứng
định lượng.

\section{Tóm tắt chương}

Chương 3 trình bày dữ liệu, cách hiệu chỉnh GBM, Heston và SBTS, hai hợp đồng,
kiến trúc mạng và quy trình MSE--CVaR. Thiết kế gồm 180 lượt huấn luyện và 360
checkpoint được đánh giá tại ba giai đoạn thị trường. So sánh thống kê sử dụng
kiểm định ghép cặp, điều chỉnh Benjamini--Hochberg và MCS.
